%% Annexe D : Textes juridiques et normatifs
\chapter{Textes Juridiques et Normatifs}
\label{ann:textes}

Seuls les éléments vérifiés dans le texte officiel sont cités dans les
chapitres. Ce tableau indique où les trouver.

\section*{Droit camerounais et africain}
\begin{tableaucours}{L{5.0cm}L{4.6cm}L{2.6cm}}{%
    \tableauHeader{Texte} & \tableauHeader{Objet} & \tableauHeader{Chapitre}
}
Loi n° 2010/012 du 21 déc. 2010 & Cybersécurité et cybercriminalité & \ref{chap:enjeux} \\
Loi n° 2024/017 du 23 déc. 2024 & Protection des données personnelles & \ref{chap:conformite} \\
Convention de Malabo (UA, 2014) & Cybersécurité et données, Afrique & \ref{chap:conformite} \\
\end{tableaucours}

\section*{Normes et référentiels internationaux}
\begin{tableaucours}{L{4.4cm}L{5.2cm}L{2.6cm}}{%
    \tableauHeader{Référence} & \tableauHeader{Objet} & \tableauHeader{Chapitre}
}
ISO/IEC 27000 & Vocabulaire et vue d'ensemble & \ref{chap:iso27000} \\
ISO/IEC 27001 & Exigences du SMSI & \ref{chap:iso27000} \\
ISO/IEC 27002 & Mesures de sécurité & \ref{chap:iso27002} \\
ISO/IEC 27005 & Gestion des risques & \ref{chap:methodes} \\
ISO 31000 & Management du risque & \ref{chap:vocabrisque} \\
ISO 22301 & Continuité d'activité & \ref{chap:continuite} \\
ISO 19011 & Audit des systèmes de management & \ref{chap:audit} \\
NIST CSF 2.0 & Cadre de cybersécurité & \ref{chap:iso27000} \\
EBIOS Risk Manager (ANSSI) & Appréciation des risques & \ref{chap:methodes} \\
\end{tableaucours}
